% description.tex
% Project description / documentation (LaTeX)
\documentclass[11pt,a4paper]{article}
\usepackage[utf8]{inputenc}
\usepackage[T1]{fontenc}
% Use English as the document language to avoid missing babel language files on some systems.
\usepackage[english]{babel}
\usepackage{geometry}
\geometry{margin=1in}
\usepackage{hyperref}

\title{Morphing — Description}
\author{\texttt{Irina Kalinina NE:22504319}}
\date{\today}

\begin{document}
\maketitle
% Write your project description here.

\section*{Résumé}

Outil de morphing entre deux images : l'utilisateur sélectionne des
couples de points correspondants, le programme triangule les positions
interpolées et génère des images intermédiaires qui sont assemblées en
une vidéo.

\section{Fonction \texttt{main} et affichage}

\texttt{main} prépare la structure MORPH, convertit/charge les images et
appelle \texttt{afficher}. La fonction \texttt{afficher} ouvre la fenêtre (SDL),
affiche les deux images et gère la boucle d'interaction où l'on clique
gauche/droite pour créer des couples. Deux boutons existent :
Sauvegarder (valide les couples) et Quitter (selon l'état de
saved quitte immédiatement ou ferme l'interface puis lance la
triangulation).

\section{Triangulation}

Le module assemble d'abord les points (4 coins + couples) via
			\texttt{build\_all\_points}. La triangulation est construite avec
			\texttt{triangulate\_frame} qui appelle \texttt{triangulation\_from\_points}
(insertion séquentielle, tests par aires signées). On peut choisir
de trianguler chaque frame indépendamment (méthode robuste mais plus
lente) ou bien de calculer une triangulation unique au frame médian
et de réutiliser sa connectivité pour toutes les images en
interpolant uniquement les positions des sommets via
	\texttt{create\_interpolated\_tri\_from\_fixed} (approche plus rapide qui
préserve la topologie entre frames).

\section{Images intermédiaires}

Pour chaque frame, \texttt{build\_intermediate\_image} parcourt tous les
pixels, trouve le triangle contenant le pixel (coordonnées
barycentriques), projette le point dans les images source/destination
et mélange les couleurs (nearest-neighbor + interpolation linéaire
selon le paramètre d'interpolation). Cette fonction remplit la
structure \texttt{IMAGE} et la retourne pour écriture.

\section{Génération de la vidéo}

Les images sont sauvegardées en PPM dans \texttt{images\_transformed/}.
Puis assemblées en MP4 par \texttt{ffmpeg} (règle dans le Makefile).

\end{document}
